% !TeX program = xelatex
% ============================================================================
%  第 7 课　群里的小世界：阶与子群
% ----------------------------------------------------------------------------
%  编译        XeLaTeX，编译两遍
%  样式        ../../00-共享/preamble.tex
%  说明        本课不发单页
%  记号        同第 6 课：e 不动，r 转 120°，r² 转 240°，f 沿过 A 的轴翻，
%              rf 沿过 B 的轴翻，r²f 沿过 C 的轴翻
% ============================================================================

\documentclass[10pt]{beamer}
% ============================================================================
%  讲义样式文件 —— 字体 / 配色 / 校徽栏 / 顶部导航栏 / 页脚 / 常用方框
%  所有个性化设置都集中在这里，改一次即可用于所有讲义
%
%  用法：主文件里写
%      \documentclass[10pt]{beamer}
%      \input{preamble.tex}
%  依赖：同目录下的 logo_white.png（顶部校徽）
%  编译：必须用 XeLaTeX，而且必须编译两遍
%        （顶部导航栏的节名与小圆点记录在 .nav 文件里，只编一遍不会更新）
% ============================================================================
\usepackage{ctex}  % 支持中文
\usepackage{amsmath, amssymb}  % 数学公式支持
\usepackage{fontspec}  % 字体设置
\usepackage{unicode-math}  % 支持Unicode数学字体
\usepackage{xcolor}  % 颜色
\usepackage{tikz}  % 用于绘制页脚色条
\usepackage{booktabs}  % 表格

% 图片搜索路径 ------------------------------------------------------------------
% 校本课的资料分目录存放，本文件可能被 02-课件/第XX课/ 、03-讲义/ 下的主文件调用，
% 这里把几种可能的相对位置都列上，logo_white.png 放在 00-共享/ 即可被找到。
% 最后一项 00-共享/图片/ 是给 \imgfig 用的插图目录（图与课件不在同一层，靠它找到）。
\graphicspath{{./}{./图片/}{../../00-共享/}{../00-共享/}{00-共享/}{../../00-共享/图片/}{../00-共享/图片/}{00-共享/图片/}}

% 设置字体
\setCJKmainfont{SimSun}  % 设置中文字体为宋体
\setmainfont{Palatino Linotype}  % 设置英文字体为Palatino Linotype
\setmathfont{Latin Modern Math}  % 设置数学字体为Latin Modern Math
% 注意：beamer 默认把正文字体族设成无衬线（sans），Windows 下对应的中文字体是
% 微软雅黑，所以上面这行 SimSun 目前对正文不起作用。想让正文变成宋体，取消
% 下面这行的注释即可（标题栏、页脚、方框标题不受影响）。
% \renewcommand{\familydefault}{\rmdefault}

% 定义颜色 - 均为绿色系
\definecolor{maingreen}{RGB}{34,139,34}  % 森林绿（导航栏底色）
\definecolor{lightgreen}{RGB}{240,255,240}  % 蜜瓜白（帧标题栏底色）
\definecolor{accentgreen}{RGB}{46,139,87}  % 海绿（强调色）
\definecolor{textgreen}{RGB}{0,128,0}  % 纯绿（标题文字）
\definecolor{darkgreen}{RGB}{20,90,20}  % 深绿（顶部校徽栏底色）

% 设置Beamer主题
\usetheme{default}
\usecolortheme{default}

% 自定义主题
\setbeamercolor{background canvas}{bg=white}
\setbeamercolor{frametitle}{fg=textgreen, bg=lightgreen}
\setbeamercolor{title}{fg=textgreen}
\setbeamercolor{subtitle}{fg=textgreen}
\setbeamercolor{section in toc}{fg=accentgreen}
\setbeamercolor{normal text}{fg=black}
\setbeamercolor{structure}{fg=accentgreen}
\setbeamercolor{itemize item}{fg=accentgreen}
\setbeamercolor{itemize subitem}{fg=accentgreen}
\setbeamercolor{itemize subsubitem}{fg=accentgreen}
\setbeamercolor{block title}{fg=white, bg=accentgreen}
\setbeamercolor{block body}{bg=lightgreen}
% example 环境（底层是 exampleblock）用的是独立的一组颜色名，default 配色主题下
% 它们的底色是空的，所以只设 block 的话“例”会变成一行没有框的标题，这里接过来
\setbeamercolor{block title example}{parent=block title}
\setbeamercolor{block body example}{parent=block body}

% ==================== 常用定理类方框（五类配色）====================
% 五类各自的标题条底色，正文底色由同色自动调浅（!8!white）
\definecolor{boxThm}{RGB}{20,90,20}    % 一、定理类：深绿
\definecolor{boxDef}{RGB}{46,139,87}   % 二、定义类：海绿
\definecolor{boxExm}{RGB}{34,139,34}   % 三、例题类：森林绿
\definecolor{boxPrf}{RGB}{85,107,47}   % 四、过程类：橄榄绿
\definecolor{boxNote}{RGB}{0,121,107}  % 五、注记类：青绿

% 方框本体：先套本类配色，再调用 beamer 的 block（改动只在环境内部生效）
\newcommand{\gbox}[2]{%
  \setbeamercolor{block title}{fg=white,bg=#1}%
  \setbeamercolor{block body}{bg=#1!8!white}%
  \begin{block}{#2}}
\newcommand{\gboxend}{\end{block}}
% 标题后的可选补充名：\begin{theorem}[勾股定理] → 定理（勾股定理）
\newcommand{\gsub}[1]{\ifx\relax#1\relax\else（#1）\fi}


% 一、定理类（定理/推论/性质/命题/引理）
\renewenvironment{theorem}[1][]{\gbox{boxThm}{定理\gsub{#1}}}{\gboxend}
\renewenvironment{corollary}[1][]{\gbox{boxThm}{推论\gsub{#1}}}{\gboxend}
\newenvironment{character}[1][]{\gbox{boxThm}{性质\gsub{#1}}}{\gboxend}
\newenvironment{proposition}[1][]{\gbox{boxThm}{命题\gsub{#1}}}{\gboxend}
\renewenvironment{lemma}[1][]{\gbox{boxThm}{引理\gsub{#1}}}{\gboxend}

% 二、定义类（定义/公理）
\renewenvironment{definition}[1][]{\gbox{boxDef}{定义\gsub{#1}}}{\gboxend}
\newenvironment{axiom}[1][]{\gbox{boxDef}{公理\gsub{#1}}}{\gboxend}

% 三、例题类（例/练习/习题）
\renewenvironment{example}[1][]{\gbox{boxExm}{例\gsub{#1}}}{\gboxend}
\newenvironment{exercise}[1][]{\gbox{boxExm}{练习\gsub{#1}}}{\gboxend}
\newenvironment{homework}[1][]{\gbox{boxExm}{习题\gsub{#1}}}{\gboxend}

% 四、过程类（证明/解/法1~法4）
\renewenvironment{proof}[1][]{\gbox{boxPrf}{证明\gsub{#1}}}{\qedsymbol\gboxend}
\newenvironment{jie}[1][]{\gbox{boxPrf}{解\gsub{#1}}}{\gboxend}
\newenvironment{way1}[1][]{\gbox{boxPrf}{法1\gsub{#1}}}{\gboxend}
\newenvironment{way2}[1][]{\gbox{boxPrf}{法2\gsub{#1}}}{\gboxend}
\newenvironment{way3}[1][]{\gbox{boxPrf}{法3\gsub{#1}}}{\gboxend}
\newenvironment{way4}[1][]{\gbox{boxPrf}{法4\gsub{#1}}}{\gboxend}

% 五、注记类（经验总结/注/值得注意的是/分析/提示）
\newenvironment{experience}[1][]{\gbox{boxNote}{经验总结\gsub{#1}}}{\gboxend}
\newenvironment{zhu}[1][]{\gbox{boxNote}{注\gsub{#1}}}{\gboxend}
\newenvironment{att}[1][]{\gbox{boxNote}{值得注意的是\gsub{#1}}}{\gboxend}
\newenvironment{analyse}[1][]{\gbox{boxNote}{分析\gsub{#1}}}{\gboxend}
\newenvironment{hint}[1][]{\gbox{boxNote}{提示\gsub{#1}}}{\gboxend}

% 顶部导航栏（miniframes 主题：显示各节名称 + 帧圆点）
% 节名与圆点记录在 .nav 文件中，改动章节后需要编译两遍才会同步
\useoutertheme[subsection=false]{miniframes}
\setbeamercolor{section in head/foot}{fg=yellow, bg=maingreen}
\setbeamercolor{mini frame}{fg=yellow, bg=maingreen}
\setbeamercolor{mini frame in current subsection}{fg=yellow, bg=accentgreen}

% 顶部校徽栏：位于导航栏上方，宽度与帧标题栏一致（整页宽）
% 校训用的楷体：ctex 的 \kaishu 在本机指向未安装的 Adobe 楷体，这里直接指定系统楷体
\setCJKfamilyfont{kaiti}{KaiTi}
\setbeamercolor{logo bar}{fg=white, bg=darkgreen}
\addtobeamertemplate{headline}{%
  \begin{beamercolorbox}[wd=\paperwidth,ht=0.84cm,dp=0.20cm,leftskip=0.25cm,rightskip=0.25cm]{logo bar}%
    % -0.035cm：色条内容默认坐在基线上，会比栏的中线略高，下移这一点使其上下居中
    \raisebox{-0.035cm}{\includegraphics[height=0.72cm]{logo_white.png}}%
    % 右侧校训：字体为上面定义的楷体，也可改为 \insertsectionhead（当前节名）、\inserttitle（章节名）等
    % 0.375cm = 校徽中线（0.36cm）加汉字基线偏移，并同样下移 0.035cm 与校徽对齐
    \hfill\raisebox{\dimexpr 0.375cm - 0.5\height - 0.5\depth\relax}%
      {\CJKfamily{kaiti}\fontsize{12pt}{24pt}\selectfont 自强不息，止于至善}%
  \end{beamercolorbox}%
}{}

% % 设置边框 - 使用TikZ绘制圆角
% \setbeamertemplate{frametitle}{
%   \vspace*{0.2cm}
%   \begin{tikzpicture}
%     % 背景矩形宽度限制为标题内容宽度
%     \fill[lightgreen, rounded corners=5pt] (0,0) rectangle (\textwidth+0.6cm, 0.8cm);
%     \node[anchor=west, text width=\textwidth, xshift=0.3cm, yshift=0.4cm] {\usebeamerfont{frametitle}\insertframetitle};
%   \end{tikzpicture}
% }

% 设置项目符号
\setbeamertemplate{itemize item}{\small\raise0.5pt\hbox{\textcolor{accentgreen}{$\bullet$}}}
\setbeamertemplate{itemize subitem}{\small\raise0.3pt\hbox{\textcolor{accentgreen}{$\circ$}}}

% 设置页脚 - 使用TikZ绘制一条整页宽的色条，左侧页码、右侧姓名
\setbeamertemplate{footline}{
  \begin{tikzpicture}
    \fill[maingreen] (-1,0) rectangle (\paperwidth, 0.4cm);
    \node[yellow,anchor=east, xshift=0cm, yshift=0.2cm] {\insertframenumber/\inserttotalframenumber};
    % 原点在纸张左边缘右侧 1cm 处（色条从 x=-1cm 画起），故 \paperwidth-2cm
    % 对应纸张右边缘内缩 1cm 的位置，与左侧页码左右对称
    \node[yellow,anchor=east, xshift=\paperwidth-2cm, yshift=0.2cm] {\insertauthor};
  \end{tikzpicture}
}

% 关闭 beamer 默认的导航符号（页码与姓名已在页脚中自行绘制）
\setbeamertemplate{navigation symbols}{}

% 设置标题页面
\setbeamertemplate{title page}{
  \vspace{0.5cm}
  \begin{center}
    \begin{minipage}[c][0.68\paperheight][c]{\textwidth}
      {\usebeamerfont{title}\usebeamercolor[fg]{title}\large\inserttitle\par}
      \vspace{1.5cm}
      {\usebeamerfont{subtitle}\usebeamercolor[fg]{subtitle}\Large\insertsubtitle\par}
      \vspace{1.5cm}
      {\usebeamerfont{author}\usebeamercolor[fg]{author}\insertauthor\par}
      \vspace{0.5cm}
      {\usebeamerfont{institute}\usebeamercolor[fg]{institute}\insertinstitute\par}
      \vspace{0.5cm}
      {\usebeamerfont{date}\usebeamercolor[fg]{date}\insertdate\par}
    \end{minipage}
  \end{center}
}

% 设置块环境
\setbeamertemplate{blocks}[rounded][shadow=false]

% 设置目录
\setbeamertemplate{section in toc}{\inserttocsectionnumber.~\inserttocsection}


% ----------------------------------------------------------------------------
%  图片占位符（临时用；有图以后把那行换成 \includegraphics 即可）
%      \imgph{说明文字}          默认高 3.6cm
%      \imgph[6cm]{说明文字}     指定高度
%  会画一个绿色虚线框，框里写说明，方便对着 PDF 找位置。
\newcommand{\imgph}[2][3.6cm]{%
  \begin{center}%
  \begin{tikzpicture}%
    \node[draw=accentgreen, densely dashed, line width=1pt, rounded corners=4pt,
          fill=lightgreen, inner sep=8pt, align=center, text width=0.72\textwidth,
          minimum height=#1, font=\small]
      {\textcolor{accentgreen}{\bfseries 【图片占位】}\\[4pt] #2};%
  \end{tikzpicture}%
  \end{center}%
}

% ----------------------------------------------------------------------------
%  插图（图已就位后用这个；图放在 00-共享\图片\ 下，由 \graphicspath 找到）
%      \imgfig{文件名}                 默认宽 0.72\textwidth
%      \imgfig[0.55\textwidth]{文件名}  指定宽度
%      \imgfig[8cm]{文件名}            也可以直接给尺寸
%  图源与编译方法见 00-共享\图片\README.md。每张图是一页 PDF，只含图、无页边。
\newcommand{\imgfig}[2][0.72\textwidth]{%
  \begin{center}%
    \includegraphics[width=#1]{#2}%
  \end{center}%
}

% ----------------------------------------------------------------------------
%  可用方框环境速查（五类配色，均不编号，均支持 \begin{xxx}[补充名]）
%    定理类（深绿）  theorem 定理 / corollary 推论 / character 性质 /
%                    proposition 命题 / lemma 引理
%    定义类（海绿）  definition 定义 / axiom 公理
%    例题类（森林绿）example 例 / exercise 练习 / homework 习题
%    过程类（橄榄绿）proof 证明 / jie 解 / way1 法1 / way2 法2 /
%                    way3 法3 / way4 法4
%    注记类（青绿）  experience 经验总结 / zhu 注 / att 值得注意的是 /
%                    analyse 分析 / hint 提示
%  另有 beamer 自带的普通方框：\begin{block}{标题} ... \end{block}
% ----------------------------------------------------------------------------


% ==================== 每份课件只改这一段 ====================
\title{从数数开始的抽象代数生活}
\subtitle{第 7 课　群里的小世界：阶与子群}
\author{姜博瀚}
\institute{厦门大学附属实验中学}
\date{\today}
% ============================================================

\begin{document}

\begin{frame}
  \titlepage
\end{frame}

% ============================================================================
\section{回顾}

\begin{frame}{回顾：乘法表里那两行}
  上节课的乘法表，\(r\) 行和 \(f\) 行都是：
  \[
    r,\ r^2,\ e,\ rf,\ r^2f,\ f \quad(\text{换个次序}) .
  \]
  \pause
  \begin{center}
    每一行都是同样 \(6\) 个动作的\textbf{一次重排}。
  \end{center}
  \pause
  \begin{center}
    今天在小范围里再看一遍这件事：一个动作自己反复做，会怎么样？
  \end{center}
\end{frame}

% ============================================================================
\section{元素的阶}

\begin{frame}{一个动作，连着做}
  \begin{definition}[元素的阶]
    一个动作 \(g\) 连续做 \(n\) 次，\textbf{第一次}回到「不动」，这个 \(n\)
    就叫做 \(g\) 的\textbf{阶}。
  \end{definition}
  \pause
  \begin{example}
    \(r\) 是「转 \(120^\circ\)」，连做 \(3\) 次：\(r\)，\(r^2\)，\(r^3=e\)。
    所以 \(r\) 的阶是 \(3\)。
  \end{example}
\end{frame}

\begin{frame}{一个一个数过来}
  \begin{center}
  \begin{tabular}{@{}lccc@{}}
    \toprule
    动作 & 连做 \(1\) 次 & 连做 \(2\) 次 & 连做 \(3\) 次 \\
    \midrule
    \(e\)（不动） & \(e\) & \(e\) & \(e\) \\
    \(r\)（转 \(120^\circ\)） & \(r\) & \(r^2\) & \(\mathbf{e}\) \\
    \(f\)（沿过 \(A\) 的轴翻） & \(f\) & \(\mathbf{e}\) & \(\mathbf{f}\) \\
    \bottomrule
  \end{tabular}
  \end{center}
  \pause
  \begin{center}
    {\large 所以：\(e\) 的阶是 \(1\)，\(r\) 的阶是 \(3\)，\(f\) 的阶是 \(2\)。}
  \end{center}
\end{frame}

\begin{frame}{阶 \(1\)、\(2\)、\(3\)，是什么意思}
  \begin{itemize}
    \item 阶 \(1\)：做一次就回到原位（只有「不动」）。
    \item 阶 \(2\)：做两次回到原位（三个翻折都是）。
    \item 阶 \(3\)：做三次回到原位（两个转动都是）。
  \end{itemize}
  \pause
  \begin{att}
    「阶」就是「这个动作要重复几次，才转完一圈回到自己」。
  \end{att}
\end{frame}

\begin{frame}{把两行式写短一点}
  转 \(120^\circ\) 那张两行式，可以缩写成一个圈：
  \[
    \begin{pmatrix}A&B&C\\B&C&A\end{pmatrix}
    \quad\rightsquigarrow\quad
    (A\ B\ C)
  \]
  \pause
  读法：\(A\) 去 \(B\)，\(B\) 去 \(C\)，\(C\) 回到 \(A\)。
  \pause
  \begin{att}
    圈里有几个字母，阶就是几。\((A\ B\ C)\) 里有 \(3\) 个字母，阶就是 \(3\)。
    这个写法只是省事，今天看一眼就行。
  \end{att}
\end{frame}

\begin{frame}{用手指回答}
  \begin{itemize}
    \item \(r^2\)（转 \(240^\circ\)）的阶是多少？
    \item \(rf\)（沿过 \(B\) 的轴翻）的阶是多少？
    \item \((A\ B)\) 这个圈的阶是多少？
  \end{itemize}
  \pause
  \begin{block}{答案}
    \(3\)；\quad \(2\)；\quad \(2\)。
  \end{block}
\end{frame}

% ============================================================================
\section{子群}

\begin{frame}{群里面能不能再装一个群}
  \begin{definition}[子群]
    设 \(H\) 是群 \(G\) 的一部分。如果 \(H\) 按同一个运算\textbf{自己也构成一个群}，
    就说 \(H\) 是 \(G\) 的一个\textbf{子群}。
  \end{definition}
  \pause
  \begin{center}
    怎么检查？只看一条最要紧的：\textbf{会不会跑出去}。
  \end{center}
\end{frame}

\begin{frame}{先看三个转动的圈子}
  圈出 \(H=\{e,\ r,\ r^2\}\)：三个转动的动作。
  \pause
  \[
  \begin{array}{@{}c|ccc@{}}
    \toprule
    & e & r & r^2 \\
    \midrule
    e & e & r & r^2 \\
    r & r & r^2 & e \\
    r^2 & r^2 & e & r \\
    \bottomrule
  \end{array}
  \]
  \pause
  \begin{center}
    表格里的每一个结果都还在这三个里——\textbf{没有跑出去}。\\[0.4em]
    所以 \(\{e,r,r^2\}\) 是一个子群。
  \end{center}
\end{frame}

\begin{frame}{再看「不动加一次翻折」}
  圈出 \(H=\{e,\ f\}\)：
  \[
  \begin{array}{@{}c|cc@{}}
    \toprule
    & e & f \\
    \midrule
    e & e & f \\
    f & f & e \\
    \bottomrule
  \end{array}
  \]
  \pause
  \begin{center}
    \(f\) 连做两次就回到 \(e\)，也没有跑出去。所以 \(\{e,f\}\) 也是子群。
  \end{center}
\end{frame}

\begin{frame}{试一个不成的}
  圈出 \(H=\{e,\ r,\ f\}\)，三个元素。\pause 检查表格：
  \[
  \begin{array}{@{}c|ccc@{}}
    \toprule
    & e & r & f \\
    \midrule
    e & e & r & f \\
    r & r & r^2 & \mathbf{r^2f} \\
    f & f & \mathbf{rf} & e \\
    \bottomrule
  \end{array}
  \]
  \pause
  \begin{center}
    \(r\cdot r=r^2\) 不在里面，\(r\cdot f=fr=r^2f\) 也不在里面。\\[0.4em]
    \textbf{跑出去了}，所以 \(\{e,r,f\}\) 不是子群。
  \end{center}
  \pause
  \begin{center}
    {\small 看仔细：\(r\cdot f=r^2f\)，可 \(f\cdot r=rf\)——两个结果不一样。\\
    这和第 6 课那句 \(rf\neq fr\) 是同一件事。}
  \end{center}
\end{frame}

\begin{frame}{\(S_3\) 里的全部子群}
  \begin{center}
  \begin{tabular}{@{}llc@{}}
    \toprule
    子群 & 是谁 & 大小 \\
    \midrule
    \(\{e\}\) & 只有不动 & \(1\) \\
    \(\{e,\ r,\ r^2\}\) & 三个转动 & \(3\) \\
    \(\{e,\ f\}\) & 不动加一条轴的翻折 & \(2\) \\
    \(\{e,\ rf\}\) & 同上 & \(2\) \\
    \(\{e,\ r^2f\}\) & 同上 & \(2\) \\
    \(S_3\) & 全部六个 & \(6\) \\
    \bottomrule
  \end{tabular}
  \end{center}
  \pause
  \begin{center}
    {\large 小群的大小 \(1,2,3,6\)，每一个都能整除大群的大小 \(6\)。}
  \end{center}
\end{frame}

\begin{frame}{这节课要带走的两句话}
  \begin{itemize}
    \item 元素的\textbf{阶}：这个动作重复几次回到原位。
    \item 群里面还能装\textbf{小群}（子群）——检查的办法是看它会不会跑出去。
  \end{itemize}
  \pause
  \begin{center}
    {\large\textcolor{textgreen}{\bfseries 下节课：换一张更大的纸——正方形。}}
  \end{center}
\end{frame}

% ---------------- 结尾页 ----------------
\begin{frame}
  \begin{center}
    \vspace{1.6cm}
    {\Large\textcolor{textgreen}{\bfseries 下课}}

    \vspace{1.0cm}
    {\large\textcolor{accentgreen}{下节课：正方形的对称——\(D_4\) 的 \(8\) 个元素}}
  \end{center}
\end{frame}

\end{document}